\section{Tame weights and integral character comparison}\label{sec:formal-comparison}
We now prove a uniform comparison theorem for two supersingular groups. The point of the argument is to estimate multiplication by $p$ on the quotient appearing in the graph construction. Its coefficients have extra divisibility forced by the common determinant.

\begin{lemma}\label{lem:weighted-kernel}
Let $p\ge5$, $N=p^2-1$, and let $L/K$ be cyclic, totally tamely ramified of degree $N$, with algebraically closed residue field. Write $R=\Ocal_L$, $a=v_L(p)$, and choose $\tau\pi=\zeta\pi$ with $\zeta\in K$ of order $N$. Let $G/R$ be a formal group of finite height with semilinear descent action. Suppose that its special fibre is a product of one-dimensional groups of height two, and that the exponents $c_i$ of the cotangent action satisfy
\[
                         c_i\equiv c_j\pmod{p-1}.
\]
Define $q_0=1$ and
\begin{equation}\label{eq:weighted-recurrence}
                   q_{j+1}=\min\{a+q_j,pq_j+p-1\}.
\end{equation}
If a homomorphism $f:F\to G$ has finite special kernel and no nonzero finite flat subgroup in its kernel, then
\[
 q_u\ge a/(p-1)\quad\Longrightarrow\quad F^u\ker(f_k)=0,
\]
where $F^u$ denotes relative Frobenius iterated $u$ times.
\end{lemma}
\begin{proof}
Average lifts of a cotangent eigenbasis by the tame idempotents. The resulting parameters $X_i$ satisfy $\tau X_i=\zeta^{c_i}X_i$ exactly. If $b_{j,\alpha}X^\alpha$ occurs in the $j$th component of $[p]$, equivariance gives
\[
       \tau b_{j,\alpha}=\zeta^{c_j-\sum_i\alpha_i c_i}b_{j,\alpha}.
\]
The valuation of a nonzero eigenvector of $\tau$ is congruent to its exponent modulo $N$: divide by that power of $\pi$ and reduce the unit. For $|\alpha|=p$ the displayed exponent is divisible by $p-1$. The order of the reduced multiplication map is $p^2$, in these coordinates as in any coordinates with invertible linear part. Every nonzero coefficient of total degree $p$ consequently has valuation at least $p-1$.

Let $U(X)dX$ be a normalized invariant differential matrix. The identity
\[
                 U([p]X)J_{[p]}(X)=pU(X)
\]
shows that every entry of the Jacobian is divisible by $p$. A coefficient of a monomial having some exponent prime to $p$ is therefore divisible by $p$. The linear part is $pX$; all coefficients of degrees less than $p^2$ vanish on reduction. Substitution of a vector in $\pi^q tR[[t]]^d$, for $q\ge1$, now has coefficient valuations at least
\[
 \min\{a+q,pq+p-1,2pq+1,p^2q\}
                      =\min\{a+q,pq+p-1\}.
\]
Indeed the first term handles the linear part and monomials with an exponent prime to $p$, the second handles total degree $p$, the third the remaining degrees below $p^2$, and the fourth all later terms. Iteration proves the bound $q_u$ for $[p]^u(\pi h)$.

We give the integral passage from this bound to the kernel, because a bound for torsion points would not suffice. A normalized logarithm has a coefficient of total degree $r$ of valuation at least $-a\lfloor\log_p r\rfloor$. This follows by integrating the integral invariant differentials and choosing a nonzero exponent in the monomial. Since $\lfloor\log_p r\rfloor\le(r-1)/(p-1)$, substitution into the logarithm at coefficient valuation $q_u\ge a/(p-1)$ yields
\begin{equation}\label{eq:weighted-logarithm}
              p^u\log_G(\pi h)=\log_G([p]^u(\pi h))
                              \in\pi^{q_u}R[[t]]^d.
\end{equation}
This is a formal, coefficientwise assertion, so equality at the bound is allowed.

Here are the Cartier-module facts used to apply it. Put $\mathcal R=R[[\Delta]][1/p]$ and $J=\pi^{\lceil a/(p-1)\rceil}R$. Write $L\{\{\Delta\}\}$ for two-sided series with coefficient valuations bounded below and tending to infinity as the exponent tends to $-\infty$; $J\{\{\Delta\}\}$ imposes the additional condition that every coefficient lies in $J$. The skew series ring $\Omega$ has coefficients in $W(k)$ with this decay condition and product
\[
 (a_i\Delta^i)(b_j\Delta^j)=\sigma^j(a_i)b_j\Delta^{i+j},
\]
where $\sigma$ is Witt Frobenius. Set $G_\pi(X,Y)=\pi^{-1}G(\pi X,\pi Y)$, and let $E(G_k)$ denote the covariant Dieudonn\'e module. The invariant Cartier module $D_G\subset L[[\Delta]]^d$ is obtained from logarithms of integral $p$-typical curves. Its completion $\mathcal M_G=\Omega D_G\subset L\{\{\Delta\}\}^d$ has
\[
 \mathcal M_G=\Omega\pi D_{G_\pi},\qquad
 J[[\Delta]]^d\subset D_G,
 \qquad E(G_k)=D_G/(D_G\cap\mathcal R^d).
\]
For every matrix $A$ over $L$, the completion embedding is the identity on bounded elements of $AD_G$, and a vector maps to a bounded element only if it was bounded. These are \citet[Proposition~9 and Lemma~3, p.~674; Proposition~10(4), p.~675; Proposition~11(1)--(2), pp.~675--676]{Bondarko2006}. The ring $\Omega$ has integral Witt coefficients and preserves coefficient valuation ideals. Equation~\eqref{eq:weighted-logarithm} therefore implies
\[
                         p^u\mathcal M_G\subset J\{\{\Delta\}\}^d.
\]
Consequently, if $A$ is integral, $z\in D_G$, and $Az$ is bounded, then $p^uAz$ has every coefficient in $J$. Since
\[
 \mathbf f^u\left(\sum_{j\ge0}z_j\Delta^j\right)
                  =p^u\sum_{j\ge0}z_{j+u}\Delta^j,
\]
we obtain $\mathbf f^uAz\in J[[\Delta]]$.

Put $I'=df(0)L^{\dim F}$. The absence of a finite flat kernel gives
\[
                  df(0)D_F=D_G\cap I'[[\Delta]].
\]
To see this, the right side corresponds to a formal group by the Cartier classification; $F$ maps to it by an isogeny, which must be an isomorphism because its kernel would be finite flat. Choose an integral complement to $I'\cap R^d$ and let $\operatorname{pr}$ be the resulting projection. If the reduction of $z\in D_G$ lies in the rational saturation of $E(F_k)$ inside $E(G_k)$, then $\operatorname{pr}z$ is bounded. The previous paragraph gives
\[
 c=\mathbf f^u\operatorname{pr}z\in J[[\Delta]]^d\subset D_G.
\]
It reduces to zero, while $\mathbf f^uz-c$ lies in $df(0)D_F$. Thus $\mathbf f^u$ kills the saturation quotient. By the covariant kernel identity, this quotient is the Dieudonn\'e module of $\ker(f_k)$ and $\mathbf f$ is relative Frobenius. This last argument is the kernel argument of \citet[Proposition~16 and proof, p.~681; Theorem~2(8), equation~(6), p.~668]{Bondarko2006}, with the improved coefficient estimate just established. It proves the assertion.
\end{proof}

\begin{theorem}\label{thm:common-determinant}
Let $p\ge5$ and let $K$ have algebraically closed residue field and absolute index $e$. Let $M_1,M_2$ be height-two Barsotti--Tate groups acquiring supersingular good reduction, with common Tate determinant $\chi_p\psi$, where $\psi$ is finite. Set $a=e(p^2-1)$ and define $q_j$ by~\eqref{eq:weighted-recurrence}. If $n\ge2$ and $q_{n-1}\ge a/(p-1)$, an isomorphism $M_1[p^n]\simeq M_2[p^n]$ implies equality of their finite special inertia characters.
\end{theorem}
\begin{proof}
A faithful finite special inertia group embeds in the units of the endomorphism algebra of a height-two formal group. This is the quaternion division algebra over $\Q_p$. It has no element of order $p$ for $p\ge5$, since $\Q_p(\mu_p)$ has degree $p-1>2$, whereas a commutative subfield has degree at most two. Its finite inertia subgroup is therefore tame and cyclic, of order dividing $N=p^2-1$. Both groups have good models over the common tame extension $L/K$ of degree $N$. As explained at the end of Section~\ref{sec:cyclic}, one first obtains good reduction up to isogeny by crystalline descent, and then realizes the given lattice by a finite flat quotient. This justifies using the specified torsion modules on these good models.

Suppose that the two special characters differ. Their eigenvalue multisets are disjoint Frobenius orbits $\{\lambda,\lambda^p\}$ and $\{\mu,\mu^p\}$ in $\mu_N$. Integral polynomials in a tame generator project the product of the two special groups onto its factors: the corresponding characteristic polynomials have relatively prime reductions.

Let $H$ be the schematic closure of the graph of the torsion isomorphism in $\mathcal M_1[p^n]\times\mathcal M_2[p^n]$. Over $L$, the determinant is cyclotomic. The alternating Tate pairings, and their inverse duality maps, extend to the good groups by \citet[Theorem~4, p.~180, and Corollaries~1--2, p.~181]{Tate1967}. Rescaling one pairing by a unit makes $H$ isotropic for their difference. It is then its own orthogonal: both finite flat subgroups have the same rank, and their inclusion becomes equality on coordinate rings, which are free modules of that rank. Cartier duality commutes with reduction, so the same holds for $H_k$.

The tame projectors preserve $H_k$, giving $H_k=J_1\times J_2$. A finite subgroup of rank $p^r$ in a one-dimensional formal group over $k$ has ideal $(T^{p^r})$ and equals $\ker F^r$. Self-orthogonality inside each height-two $p^n$-torsion factor forces $r=n$. Hence
\[
                        H_k=\ker F_1^n\times\ker F_2^n.
\]
Form $Q=(\mathcal M_1\times\mathcal M_2)/H$ and $G=Q\times\mathcal M_1$. The map
\[
 f:\mathcal M_1\times\mathcal M_2\longrightarrow G,
             \qquad z\longmapsto(qz,\operatorname{pr}_1z)
\]
has no finite flat kernel: generically the graph intersects $0\times M_2$ trivially. Its special kernel is $0\times\ker F_2^n$.

Let $c_1,c_2$ be the two cotangent exponents. The common finite determinant is their common reduced norm; hence $c_1\equiv c_2\pmod{p-1}$. Quotient formation and relative Frobenius show that the three cotangent exponents on $G_k$ are
\[
                         p^nc_1,\quad p^nc_2,\quad c_1.
\]
They are congruent modulo $p-1$. Lemma~\ref{lem:weighted-kernel}, applied to this target with $u=n-1$, would make $F^{n-1}\ker F_2^n$ zero. Its rank is $p$, a contradiction. Therefore the characters agree.
\end{proof}
